\documentclass[10pt,a4paper]{amsart}
\usepackage{amsmath,amssymb,amsthm,graphicx}
% Retained preamble for the Stacks Project chapter snapshot in this directory.
% The chapter's own preamble.tex is a copy of the shared Stacks Project preamble; the block
% below is sliced MECHANICALLY, line for line, from it. Nothing was added, removed or
% redefined; the wrapper class, the xr-hyper/hyperref wiring and the externaldocument list
% are omitted because the assembler supplies the document class.
%
% lines taken from preamble.tex: 3, 5, 9, 11, 12, 156, 157, 158, ... (51 lines total)
%
% package coverage restored from the chapter's own preamble (2026-10-13): xy, xy-2cell

\usepackage[all]{xy}
\xyoption{2cell}
\UseAllTwocells

\newenvironment{reference}{\comment}{\endcomment}
\newenvironment{slogan}{\comment}{\endcomment}
\newenvironment{history}{\comment}{\endcomment}
\theoremstyle{plain}
\newtheorem{theorem}[subsection]{Theorem}
\newtheorem{proposition}[subsection]{Proposition}
\newtheorem{lemma}[subsection]{Lemma}
\theoremstyle{definition}
\newtheorem{definition}[subsection]{Definition}
\newtheorem{example}[subsection]{Example}
\newtheorem{exercise}[subsection]{Exercise}
\newtheorem{situation}[subsection]{Situation}
\theoremstyle{remark}
\newtheorem{remark}[subsection]{Remark}
\newtheorem{remarks}[subsection]{Remarks}
\numberwithin{equation}{subsection}
\def\lim{\mathop{\mathrm{lim}}\nolimits}
\def\colim{\mathop{\mathrm{colim}}\nolimits}
\def\Spec{\mathop{\mathrm{Spec}}}
\def\Hom{\mathop{\mathrm{Hom}}\nolimits}
\def\Ext{\mathop{\mathrm{Ext}}\nolimits}
\def\SheafHom{\mathop{\mathcal{H}\!\mathit{om}}\nolimits}
\def\SheafExt{\mathop{\mathcal{E}\!\mathit{xt}}\nolimits}
\def\Sch{\mathit{Sch}}
\def\Mor{\mathop{\mathrm{Mor}}\nolimits}
\def\Ob{\mathop{\mathrm{Ob}}\nolimits}
\def\Sh{\mathop{\mathit{Sh}}\nolimits}
\def\NL{\mathop{N\!L}\nolimits}
\def\CH{\mathop{\mathrm{CH}}\nolimits}
\def\proetale{{pro\text{-}\acute{e}tale}}
\def\etale{{\acute{e}tale}}
\def\QCoh{\mathit{QCoh}}
\def\Ker{\mathop{\mathrm{Ker}}}
\def\Im{\mathop{\mathrm{Im}}}
\def\Coker{\mathop{\mathrm{Coker}}}
\def\Coim{\mathop{\mathrm{Coim}}}
\DeclareMathSymbol{\boxtimes}{\mathbin}{AMSa}{"02}
\def\QCohstack{\mathcal{QC}\!\mathit{oh}}
\def\Cohstack{\mathcal{C}\!\mathit{oh}}
\def\Spacesstack{\mathcal{S}\!\mathit{paces}}
\def\Quotfunctor{\mathrm{Quot}}
\def\Hilbfunctor{\mathrm{Hilb}}
\def\Curvesstack{\mathcal{C}\!\mathit{urves}}
\def\Polarizedstack{\mathcal{P}\!\mathit{olarized}}
\def\Complexesstack{\mathcal{C}\!\mathit{omplexes}}
\def\Pic{\mathop{\mathrm{Pic}}\nolimits}
\def\Picardstack{\mathcal{P}\!\mathit{ic}}
\def\Picardfunctor{\mathrm{Pic}}
\def\Deformationcategory{\mathcal{D}\!\mathit{ef}}
\usepackage{cleveref}
\crefname{originallemma}{original lemma}{original lemmas}
\Crefname{originallemma}{Original lemma}{Original lemmas}
\makeatletter
\newlabel{lemma-quot-existence-part}{{0DM4}{source}{}{sourcepointer-1}{}}
\newlabel{lemma-quot-existence-part@cref}{{[originallemma][0][]0DM4}{[1][0][]source}{}{}{}}
\makeatother
\crefname{originallemma}{original lemma}{original lemmas}
\Crefname{originallemma}{Original lemma}{Original lemmas}
\makeatletter
\newlabel{lemma-quot-s-lfp}{{0DM3}{source}{}{sourcepointer-2}{}}
\newlabel{lemma-quot-s-lfp@cref}{{[originallemma][0][]0DM3}{[1][0][]source}{}{}{}}
\makeatother
\crefname{originallemma}{original lemma}{original lemmas}
\Crefname{originallemma}{Original lemma}{Original lemmas}
\makeatletter
\newlabel{more-morphisms-lemma-generic-flatness-stratification}{{0ASY}{source}{}{sourcepointer-3}{}}
\newlabel{more-morphisms-lemma-generic-flatness-stratification@cref}{{[originallemma][0][]0ASY}{[1][0][]source}{}{}{}}
\makeatother
\crefname{originallemma}{original lemma}{original lemmas}
\Crefname{originallemma}{Original lemma}{Original lemmas}
\makeatletter
\newlabel{spaces-morphisms-lemma-characterize-proper}{{0A40}{source}{}{sourcepointer-4}{}}
\newlabel{spaces-morphisms-lemma-characterize-proper@cref}{{[originallemma][0][]0A40}{[1][0][]source}{}{}{}}
\makeatother
\crefname{originallemma}{original lemma}{original lemmas}
\Crefname{originallemma}{Original lemma}{Original lemmas}
\makeatletter
\newlabel{varieties-lemma-bound-quotients-free}{{08AG}{source}{}{sourcepointer-5}{}}
\newlabel{varieties-lemma-bound-quotients-free@cref}{{[originallemma][0][]08AG}{[1][0][]source}{}{}{}}
\makeatother
\crefname{originallemma}{original lemma}{original lemmas}
\Crefname{originallemma}{Original lemma}{Original lemmas}
\makeatletter
\newlabel{varieties-lemma-m-regular-globally-generated}{{08A8}{source}{}{sourcepointer-6}{}}
\newlabel{varieties-lemma-m-regular-globally-generated@cref}{{[originallemma][0][]08A8}{[1][0][]source}{}{}{}}
\makeatother
\crefname{originallemma}{original lemma}{original lemmas}
\Crefname{originallemma}{Original lemma}{Original lemmas}
\makeatletter
\newlabel{varieties-lemma-m-regular-up}{{08A6}{source}{}{sourcepointer-7}{}}
\newlabel{varieties-lemma-m-regular-up@cref}{{[originallemma][0][]08A6}{[1][0][]source}{}{}{}}
\makeatother
\title{Quot of O\textasciicircum{}\{oplus r\} on P\textasciicircum{}n\_Z with Hilbert polynomial P is proper over Z}
\author{The Stacks Project Authors}
\date{stacks-project:moduli.tex:local-snapshot-20261012; GNU FDL 1.2 or later}
\begin{document}
\maketitle
\noindent\textit{Editorial note.} Statement span is lines 1046-1056 of sources/stacks-moduli/pinned-original-source.tex; proof span is lines 1058-1190, the complete author proof (131 body lines).\par\medskip
\noindent\textit{Editorial note.} Navigation stubs: 7 cited-but-not-included labels are supplied as non-load-bearing pointers with their pinned tags where the snapshot tag index has them.\par\medskip
\noindent\textit{Editorial note.} extra\_packages is empty on purpose: the retained preamble already loads xy with options and a second declaration clashes in the assembler wrapper.\par\medskip
\noindent\textit{Editorial note.} Text is retained verbatim, including the author's own typos and placeholders.\par\medskip
\noindent\textit{Editorial note.} Queue index 1 of queue\_moduli.json (pb=4649, stmt 1046-1056, proof 1058-1190).\par\medskip
\noindent\textit{Editorial note.} stacks-moduli wave-1: fresh adjudication accept (2026-10-14); pool per Chief order 08:5x (dpa superseded to lane-s9); moduli.tex 53 labels/48 with proof judged 25 accept / 23 reject / 5 pb=0 env-dup rows excluded; spans cite-free and reference/slogan/history-free so no patchenv needed.\par\medskip
\noindent\textit{Allowed standard prerequisites (editorial).}\begin{itemize}
\item Explicit prerequisite boundary. The cited inputs are prerequisites; the argument in the retained span is not replaced by any cited result.
\end{itemize}\medskip
\ifdefined\hypertarget\hypertarget{sourcepointer-1}{}\fi
\noindent\textit{Source reference.} Original lemma 0DM4: sources/stacks-moduli/pinned-original-source.tex, source lines 770--776. Cited by the retained statement or proof but not part of the collected spans; supplied only as a navigation pointer. The argument does not depend on its content.\par\medskip
\ifdefined\hypertarget\hypertarget{sourcepointer-2}{}\fi
\noindent\textit{Source reference.} Original lemma 0DM3: sources/stacks-moduli/pinned-original-source.tex, source lines 755--759. Cited by the retained statement or proof but not part of the collected spans; supplied only as a navigation pointer. The argument does not depend on its content.\par\medskip
\ifdefined\hypertarget\hypertarget{sourcepointer-3}{}\fi
\noindent\textit{Source reference.} Original lemma 0ASY: sources/stacks-moduli/pinned-original-source.tex, source lines 1190--1190. Cited by the retained statement or proof but not part of the collected spans; supplied only as a navigation pointer. The argument does not depend on its content.\par\medskip
\ifdefined\hypertarget\hypertarget{sourcepointer-4}{}\fi
\noindent\textit{Source reference.} Original lemma 0A40: sources/stacks-moduli/pinned-original-source.tex, source lines 1190--1190. Cited by the retained statement or proof but not part of the collected spans; supplied only as a navigation pointer. The argument does not depend on its content.\par\medskip
\ifdefined\hypertarget\hypertarget{sourcepointer-5}{}\fi
\noindent\textit{Source reference.} Original lemma 08AG: sources/stacks-moduli/pinned-original-source.tex, source lines 1190--1190. Cited by the retained statement or proof but not part of the collected spans; supplied only as a navigation pointer. The argument does not depend on its content.\par\medskip
\ifdefined\hypertarget\hypertarget{sourcepointer-6}{}\fi
\noindent\textit{Source reference.} Original lemma 08A8: sources/stacks-moduli/pinned-original-source.tex, source lines 1190--1190. Cited by the retained statement or proof but not part of the collected spans; supplied only as a navigation pointer. The argument does not depend on its content.\par\medskip
\ifdefined\hypertarget\hypertarget{sourcepointer-7}{}\fi
\noindent\textit{Source reference.} Original lemma 08A6: sources/stacks-moduli/pinned-original-source.tex, source lines 1190--1190. Cited by the retained statement or proof but not part of the collected spans; supplied only as a navigation pointer. The argument does not depend on its content.\par\medskip
\begin{lemma}
\label{lemma-quot-Pn}
Let $n \geq 0$, $r \geq 1$, $P \in \mathbf{Q}[t]$.
The algebraic space
$$
X = \Quotfunctor^P_{\mathcal{O}^{\oplus r}_{\mathbf{P}^n_\mathbf{Z}}/
\mathbf{P}^n_\mathbf{Z}/\mathbf{Z}}
$$
parametrizing quotients of $\mathcal{O}_{\mathbf{P}^n_\mathbf{Z}}^{\oplus r}$
with Hilbert polynomial $P$ is proper over $\Spec(\mathbf{Z})$.
\end{lemma}

\begin{proof}
We already know that $X \to \Spec(\mathbf{Z})$ is separated and
locally of finite presentation (Lemma \ref{lemma-quot-s-lfp}).
We also know that $X \to \Spec(\mathbf{Z})$ satisfies the
existence part of the valuative criterion, see
Lemma \ref{lemma-quot-existence-part}.
By the valuative criterion for properness, it suffices to
prove our Quot space is quasi-compact, see
Morphisms of Spaces,
Lemma \ref{spaces-morphisms-lemma-characterize-proper}.
Thus it suffices to find a quasi-compact scheme $T$ and a surjective
morphism $T \to X$. Let $m$ be the integer found in
Varieties, Lemma \ref{varieties-lemma-bound-quotients-free}.
Let
$$
N = r{m + n \choose n} - P(m)
$$
We will write $\mathbf{P}^n$ for
$\mathbf{P}^n_\mathbf{Z} = \text{Proj}(\mathbf{Z}[T_0, \ldots, T_n])$
and unadorned products will mean products over $\Spec(\mathbf{Z})$.
The idea of the proof is to construct a ``universal'' map
$$
\Psi :
\mathcal{O}_{T \times \mathbf{P}^n}(-m)^{\oplus N}
\longrightarrow
\mathcal{O}_{T \times \mathbf{P}^n}^{\oplus r}
$$
over an affine scheme $T$ and show that every point of $X$
corresponds to a cokernel of this in some point of $T$.

\medskip\noindent
Definition of $T$ and $\Psi$. We take $T = \Spec(A)$ where
$$
A = \mathbf{Z}[a_{i, j, E}]
$$
where $i \in \{1, \ldots, r\}$, $j \in \{1, \ldots, N\}$
and $E = (e_0, \ldots, e_n)$ runs through the multi-indices
of total degree $|E| = \sum_{k = 0, \ldots n} e_k = m$.
Then we define $\Psi$ to be the map whose $(i, j)$ matrix
entry is the map
$$
\sum\nolimits_{E = (e_0, \ldots, e_n)}
a_{i, j, E} T_0^{e_0} \ldots T_n^{e_n} :
\mathcal{O}_{T \times \mathbf{P}^n}(-m)
\longrightarrow
\mathcal{O}_{T \times \mathbf{P}^n}
$$
where the sum is over $E$ as above (but $i$ and $j$ are fixed of course).

\medskip\noindent
Consider the quotient $\mathcal{Q} = \Coker(\Psi)$ on $T \times \mathbf{P}^n$.
By More on Morphisms, Lemma
\ref{more-morphisms-lemma-generic-flatness-stratification}
there exists a $t \geq 0$ and closed subschemes
$$
T = T_0 \supset T_1 \supset \ldots \supset T_t = \emptyset
$$
such that the pullback $\mathcal{Q}_p$ of $\mathcal{Q}$ to
$(T_p \setminus T_{p + 1}) \times \mathbf{P}^n$ is flat over
$T_p \setminus T_{p + 1}$. Observe that we
have an exact sequence
$$
\mathcal{O}_{(T_p \setminus T_{p + 1}) \times \mathbf{P}^n}(-m)^{\oplus N}
\to
\mathcal{O}_{(T_p \setminus T_{p + 1}) \times \mathbf{P}^n}^{\oplus r}
\to
\mathcal{Q}_p
\to
0
$$
by pulling back the exact sequence defining $\mathcal{Q} = \Coker(\Psi)$.
Therefore we obtain a morphism
$$
\coprod (T_p \setminus T_{p + 1})
\longrightarrow
\Quotfunctor_{\mathcal{O}^{\oplus r}/\mathbf{P}/\mathbf{Z}}
\supset
\Quotfunctor^P_{\mathcal{O}^{\oplus r}/\mathbf{P}/\mathbf{Z}} = X
$$
Since the left hand side is a Noetherian scheme and the inclusion
on the right hand side is open, it
suffices to show that any point of $X$ is in the image of this morphism.

\medskip\noindent
Let $k$ be a field and let $x \in X(k)$. Then $x$ corresponds to
a surjection $\mathcal{O}_{\mathbf{P}^n_k}^{\oplus r} \to \mathcal{F}$
of coherent $\mathcal{O}_{\mathbf{P}^n_k}$-modules
such that the Hilbert polynomial of $\mathcal{F}$ is $P$.
Consider the short exact sequence
$$
0 \to \mathcal{K} \to
\mathcal{O}_{\mathbf{P}^n_k}^{\oplus r} \to
\mathcal{F} \to 0
$$
By Varieties, Lemma \ref{varieties-lemma-bound-quotients-free}
and our choice of $m$ we see that $\mathcal{K}$ is $m$-regular.
By Varieties, Lemma \ref{varieties-lemma-m-regular-globally-generated}
we see that $\mathcal{K}(m)$ is globally generated.
By Varieties, Lemma \ref{varieties-lemma-m-regular-up}
and the definition of $m$-regularity we see that
$H^i(\mathbf{P}^n_k, \mathcal{K}(m)) = 0$ for $i > 0$.
Hence we see that
$$
\dim_k H^0(\mathbf{P}^n_k, \mathcal{K}(m)) =
\chi(\mathcal{K}(m)) =
\chi(\mathcal{O}_{\mathbf{P}^n_k}(m)^{\oplus r}) -
\chi(\mathcal{F}(m)) = N
$$
by our choice of $N$. This gives a surjection
$$
\mathcal{O}_{\mathbf{P}^n_k}^{\oplus N}
\longrightarrow
\mathcal{K}(m)
$$
Twisting back down and using the short exact sequence above
we see that $\mathcal{F}$ is the cokernel of a map
$$
\Psi_x :
\mathcal{O}_{\mathbf{P}^n_k}(-m)^{\oplus N}
\to
\mathcal{O}_{\mathbf{P}^n_k}^{\oplus r}
$$
There is a unique ring map $\tau : A \to k$ such that the base change
of $\Psi$ by the corresponding morphism $t = \Spec(\tau) : \Spec(k) \to T$
is $\Psi_x$. This is true because the entries of the $N \times r$
matrix defining $\Psi_x$ are homogeneous polynomials
$\sum \lambda_{i, j, E} T_0^{e_0} \ldots T_n^{e_n}$
of degree $m$ in $T_0, \ldots, T_n$ with coefficients
$\lambda_{i, j, E} \in k$ and we can set
$\tau(a_{i, j, E}) = \lambda_{i, j, E}$.
Then $t \in T_p \setminus T_{p + 1}$ for some $p$ and
the image of $t$ under the morphism above is $x$ as desired.
\end{proof}

\end{document}
